\documentclass[11pt,a4paper]{article}
\usepackage[utf8]{inputenc}
\usepackage[vietnamese]{babel}
\usepackage{amsmath,amssymb}
\usepackage[version=4]{mhchem}
\usepackage{geometry}
\geometry{top=1.5cm,bottom=1.2cm,left=1.4cm,right=1.4cm,headheight=28pt,headsep=12pt,footskip=18pt}
\usepackage{graphicx}
\usepackage{xcolor}
\usepackage{tabularx}
\usepackage{array}
\usepackage{tikz}
\usepackage{fancyhdr}
\usepackage{enumitem}

\definecolor{hfRed}{HTML}{D32F2F}
\definecolor{hfGreen}{HTML}{8BC390}
\definecolor{mainGreen}{HTML}{2E7D32}

\usepackage{ex_test}

\pagestyle{fancy}
\fancyhf{}
\renewcommand{\headrulewidth}{0.75pt}
\renewcommand{\headrule}{\hbox to\headwidth{\color{hfGreen}\leaders\hrule height \headrulewidth\hfill}}
\renewcommand{\footrulewidth}{0.75pt}
\renewcommand{\footrule}{\hbox to\headwidth{\color{hfGreen}\leaders\hrule height \footrulewidth\hfill}}

\fancyhead[L]{\small\itshape\color{hfRed}Tài liệu lưu hành nội bộ}
\fancyhead[R]{\small\color{hfRed}Thầy \textbf{TRẦN VĂN THẠNH} - 0777.470.803}

\fancyfoot[L]{\footnotesize\itshape\color{hfRed}CS1: 6/15 Nguyễn Hoàng, P. Kim Long, TP Huế \quad|\quad CS2: 24 Đặng Thái Thân, TP Huế}
\fancyfoot[R]{\small\textbf{\color{hfRed}Trang \thepage}}

\setlength{\parindent}{0pt}
\setlength{\parskip}{1.5pt}

\begin{document}

\begin{minipage}[t]{0.48\textwidth}
\textbf{SỞ GD\&ĐT THÀNH PHỐ HUẾ}\\
\textbf{TRƯỜNG THPT HAI BÀ TRƯNG}\\
\textit{ĐỀ CHÍNH THỨC}\\
\textit{(Đề thi có 02 trang)}
\end{minipage}
\hfill
\begin{minipage}[t]{0.48\textwidth}
\raggedleft
\textbf{\color{hfRed}KIỂM TRA GK1 NĂM HỌC 2026 -- 2027}\\
\textbf{MÔN: HOÁ HỌC -- LỚP 12}\\
\textit{Thời gian làm bài: 45 phút (không kể thời gian phát đề)}\\
\textbf{MÃ ĐỀ: 121}
\end{minipage}

\vspace{3pt}
\noindent
\textit{Họ và tên thí sinh:} \dotfill \textit{Lớp:} \dotfill \textit{Số báo danh:} \dotfill

\vspace{2pt}
\noindent
\textit{\footnotesize (Cho biết nguyên tử khối: H = 1; C = 12; N = 14; O = 16; Na = 23; Cl = 35,5).}

\vspace{3pt}
\noindent
\textbf{\color{mainGreen}PHẦN 1: Trắc nghiệm nhiều lựa chọn}\\
\textit{Thí sinh trả lời từ câu 1 đến câu 12. Mỗi câu hỏi thí sinh chỉ chọn một phương án.}

\Opensolutionfile{ans}[ans-hbt-hs]

\begin{ex}
Ester là hợp chất hữu cơ được tạo thành khi thay thế nhóm -OH ở nhóm carboxyl (-COOH) của carboxylic acid bằng nhóm nào sau đây?
\choice
{\True \ce{-OR'}.}
{\ce{-NH2}.}
{\ce{-CH3}.}
{\ce{-Cl}.}
\end{ex}

\begin{ex}
Trong thực tế, cá tươi thường có mùi tanh đặc trưng do hợp chất trimethylamine \ce{(CH3)3N} gây ra. Khi sơ chế cá, người ta thường dùng giấm ăn hoặc chanh tươi chà xát lên thân cá để khử mùi tanh. Cơ sở hóa học của việc làm này là
\choice
{\True acid trong chanh/giấm phản ứng với trimethylamine tạo thành muối không có mùi.}
{chanh/giấm bay hơi mang theo khí trimethylamine.}
{chanh/giấm làm đông tụ protein của cá.}
{chanh/giấm làm hòa tan hoàn toàn trimethylamine bám trên da cá.}
\end{ex}

\begin{ex}
Phân tử maltose được hợp thành từ
\choice
{hai đơn vị $\beta$-glucose.}
{một đơn vị $\beta$-glucose và một đơn vị $\alpha$-fructose.}
{\True hai đơn vị $\alpha$-glucose.}
{một đơn vị $\alpha$-glucose và một đơn vị $\beta$-fructose.}
\end{ex}

\begin{ex}
Chất nào sau đây là một amine bậc 1?
\choice
{\True \ce{C2H5NH2}.}
{\ce{(CH3)3N}.}
{\ce{(CH3)2NH}.}
{\ce{CH3-NH-C2H5}.}
\end{ex}

\begin{ex}
Carbohydrate nào sau đây thuộc loại disaccharide?
\choice
{Cellulose.}
{Glucose.}
{Fructose.}
{\True Saccharose.}
\end{ex}

\begin{ex}
Ở điều kiện thường, các amino acid là những chất tinh thể không màu, dễ tan trong nước và có nhiệt độ nóng chảy cao là do trong phân tử tồn tại chủ yếu ở dạng
\choice
{phân tử trung hòa.}
{\True ion lưỡng cực.}
{anion.}
{cation.}
\end{ex}

\begin{ex}
Xà phòng là hỗn hợp muối sodium hoặc potassium của acid béo cùng các chất phụ gia. Công thức nào sau đây đại diện cho thành phần chính của xà phòng?
\choice
{\ce{C12H25-C6H4-SO3Na}.}
{\ce{CH3[CH2]10CH2OSO3Na}.}
{\True \ce{C17H35COONa}.}
{\ce{C3H5(OH)3}.}
\end{ex}

\begin{ex}
Chất nào sau đây không có khả năng tham gia phản ứng thủy phân trong môi trường acid?
\choice
{Tinh bột.}
{Cellulose.}
{\True Glucose.}
{Saccharose.}
\end{ex}

\begin{ex}
Cho các dung dịch riêng biệt sau: glucose, saccharose, ethyl alcohol. Thuốc thử duy nhất có thể dùng để phân biệt cả 3 dung dịch trên ở nhiệt độ phòng và khi đun nóng là
\choice
{Nước bromine.}
{\True \ce{Cu(OH)2} trong môi trường kiềm.}
{Quỳ tím.}
{Dung dịch \ce{AgNO3/NH3}.}
\end{ex}

\begin{ex}
Nguyên nhân chính khiến dung dịch methylamine (\ce{CH3NH2}) làm quỳ tím hóa xanh là do
\choice
{\ce{CH3NH2} bị phân hủy tạo ra \ce{NaOH}.}
{\ce{CH3NH2} có tính acid mạnh.}
{phân tử \ce{CH3NH2} chứa nhóm \ce{-OH}.}
{\True nguyên tử N còn cặp electron chưa liên kết có khả năng nhận \ce{H+} từ nước.}
\end{ex}

\begin{ex}
Số nhóm amino (\ce{-NH2}) và nhóm carboxyl (\ce{-COOH}) có trong phân tử Glycine lần lượt là
\choice
{1 và 2.}
{\True 1 và 1.}
{2 và 2.}
{2 và 1.}
\end{ex}

\begin{ex}
Trong điều kiện thường, chất nào sau đây ở trạng thái khí và có mùi khai/tanh?
\choice
{Glutamic acid.}
{Glycine.}
{\True Methylamine.}
{Aniline.}
\end{ex}

\Closesolutionfile{ans}

\vspace{5pt}
\noindent
\textbf{\color{mainGreen}PHẦN 2: Trắc nghiệm đúng-sai}\\
\textit{Thí sinh trả lời từ câu 1 đến câu 2. Trong mỗi ý a), b), c), d) ở mỗi câu, thí sinh chọn đúng hoặc sai.}

\Opensolutionfile{ans}[ans-hbt-p2-hs]

\begin{ex}
Khi để lâu ngoài không khí, dầu mỡ động thực vật thường xảy ra hiện tượng bị ôi và có mùi hôi khó chịu. Nguyên nhân chính gây ra hiện tượng này là do
\choiceTF
{Phản ứng thủy phân chất béo tạo ra glycerol có vị đắng.}
{\True Các gốc acid béo không no bị oxi hóa chậm bởi oxygen không khí.}
{Đã xảy ra phản ứng hydrogen hóa vào gốc acid béo không no.}
{Sự nhân lên của vi khuẩn phân hủy gốc acid béo no.}
\end{ex}

\begin{ex}
Tinh bột và cellulose là hai polysaccharide phổ biến trong tự nhiên, có cấu tạo và ứng dụng rất khác nhau. Xét các phát biểu sau:
\choiceTF
{Phân tử amylose trong tinh bột được cấu tạo từ nhiều đơn vị $\alpha$-glucose liên kết với nhau qua liên kết $\alpha$-1,2-glycoside tạo thành chuỗi xoắn.}
{\True Con người không thể sử dụng cellulose (như rơm, gỗ) làm thực phẩm chính cung cấp năng lượng vì trong hệ tiêu hóa của con người không có enzyme cellulase xúc tác thủy phân cellulose.}
{Phản ứng màu giữa hồ tinh bột và dung dịch iodine xuất hiện màu xanh tím do phân tử iodine đã phản ứng với các phân tử amylose có màu xanh tím. Màu sẽ biến mất khi đun nóng, để nguội lại xuất hiện.}
{Từ 1,0 tấn mùn cưa chứa 50\% cellulose về khối lượng, qua quá trình thủy phân và lên men thu được ethanol với hiệu suất toàn bộ quá trình là 80\%, khối lượng ethanol thu được là 237,2 kg.}
\end{ex}

\Closesolutionfile{ans}

\vspace{5pt}
\noindent
\textbf{\color{mainGreen}PHẦN 3: Trắc nghiệm trả lời ngắn}\\
\textit{Thí sinh trả lời từ câu 1 đến câu 4. Viết câu trả lời vào ô tương ứng.}

\begin{ex}
Một học sinh lớp 12 cần cung cấp nhu cầu năng lượng mỗi ngày là 9690 kJ, trong đó chất béo đóng góp 20\% tổng năng lượng này. Biết rằng trung bình mỗi gram chất béo khi chuyển hóa trong cơ thể cung cấp 38 kJ năng lượng. Khối lượng chất béo mà học sinh này cần tiêu thụ trong một ngày để đảm bảo dinh dưỡng là bao nhiêu gam?
\vskip 2pt
\textbf{Trả lời:} \dotfill
\end{ex}

\begin{ex}
Tính phân tử khối (theo amu) của tripeptide Ala-Gly-Val. (Biết Alanine = 89; Glycine = 75; Valine = 117).
\vskip 2pt
\textbf{Trả lời:} \dotfill
\end{ex}

\begin{ex}
Cho 8,9 gam Alanine tác dụng hoàn toàn với dung dịch \ce{NaOH} dư. Sau phản ứng thu được $m$ gam muối khan. Giá trị của $m$ là bao nhiêu? (Cho H = 1, C = 12, N = 14, O = 16, Na = 23, Cl = 35,5).
\vskip 2pt
\textbf{Trả lời:} \dotfill
\end{ex}

\begin{ex}
Cho các chất/dung dịch sau: methylamine, aniline, glycine, glutamic acid, lysine. Có bao nhiêu dung dịch chất làm quỳ tím chuyển sang màu đỏ?
\vskip 2pt
\textbf{Trả lời:} \dotfill
\end{ex}

\vspace{5pt}
\noindent
\textbf{\color{mainGreen}PHẦN 4: Tự luận}\\
\textit{Thí sinh trình bày chi tiết lời giải từ câu 1 đến câu 3.}

\begin{ex}
Tiến hành phản ứng ester hóa điều chế dầu chuối (isoamyl acetate) bằng cách đun nóng hỗn hợp gồm 16,2 g acetic acid (\ce{CH3COOH}) và 15,2 g isoamyl alcohol (\ce{(CH3)2CHCH2CH2OH}) với xúc tác \ce{H2SO4} đặc. Sau phản ứng thu được 14,16 g isoamyl acetate. Tính hiệu suất của phản ứng điều chế trên.
\vskip 50pt
\end{ex}

\begin{ex}
\phantom{a}
\begin{enumerate}[label=\alph*)]
\item Trình bày cấu tạo của phân tử saccharose, biểu diễn công thức cấu tạo dạng mạch vòng của phân tử đó (ghi rõ số chỉ vị trí carbon mỗi gốc monosaccharide).
\item Phân tử saccharose có thể mở vòng không? Giải thích.
\end{enumerate}
\vskip 60pt
\end{ex}

\begin{ex}
Thủy phân hết $m$ gam tetrapeptide Ala-Ala-Ala-Ala thu được hỗn hợp gồm 44,5 gam Ala, 24 gam Ala-Ala và 46,2 gam Ala-Ala-Ala. Xác định giá trị $m$.
\vskip 50pt
\end{ex}

\begin{center}
\textbf{---------- HẾT ----------}
\end{center}

\end{document}
